\chapter{R\'ealisation, impl\'ementation ou exp\'erimentation}

\guided{Documentez ici ce qui a \'et\'e effectivement r\'ealis\'e. Le jury attend des preuves de mise en oeuvre, d'int\'egration, de difficult\'es rencontr\'ees et de solutions adopt\'ees. Evitez les longues captures ou les blocs de code sans commentaire analytique.}

\section{Environnement mat\'eriel et logiciel}
\todo{Pr\'esenter les plateformes, biblioth\`eques, versions logicielles, capteurs, cartes, outils de simulation, IDE, frameworks ou logiciels CAO.}

\section{Etapes d'impl\'ementation}
\todo{D\'ecrire les principales fonctionnalit\'es r\'ealis\'ees, les algorithmes importants, les choix de param\'etrage.}

\begin{algorithm}[H]
\caption{Exemple de pseudo-code pour une fonctionnalit\'e principale}
\begin{algorithmic}[1]
\State Initialiser les composants / variables
\State Acqu\'erir les entr\'ees ou charger les donn\'ees
\State Appliquer le traitement / le mod\`ele / la logique de commande
\State Produire une sortie exploitable ou une action de contr\^ole
\State Journaliser les r\'esultats utiles \`a la validation
\end{algorithmic}
\end{algorithm}

\section{Difficult\'es rencontr\'ees et solutions}
\todo{Montrer la capacit\'e d'analyse et d'adaptation de l'\'etudiant.}

\section{Livrables techniques obtenus}
\todo{Prototype, code, interface, sch\'emas, PCB, maquette, tableau de bord, mod\`ele, jeu de donn\'ees enrichi, etc.}
